% STATUS: DRAFT
\chapter{Reeh--Schlieder, Bisognano--Wichmann, and the Unruh effect}\label{ch:rs_bw}

\begin{physmotiv}
Two theorems give the vacuum of a QFT its algebraic character. Reeh--Schlieder says that acting with operators from a tiny region on the vacuum produces (approximately) \emph{every} state: the vacuum is cyclic and separating for local algebras, so modular theory applies. Bisognano--Wichmann says what the modular flow is: for a Rindler wedge, it is the Lorentz boost. Together they make the Unruh effect a theorem about modular flow and show that the abstract Tomita--Takesaki structure of \cref{ch:tomita} is, in QFT, geometry.
\end{physmotiv}

\section{Reeh--Schlieder}

\begin{theorem}[Reeh--Schlieder]
Let $O$ be a non-empty open region. Under the axioms of \cref{ch:haagkastler}, the vacuum $\Omega$ is cyclic for $\A(O)$. If moreover the causal complement $O'$ has non-empty interior, $\Omega$ is also separating for $\A(O)$ (it is cyclic for $\A(O')\subset\A(O)'$).
\end{theorem}

\begin{proofsketch}
Let $\Psi\perp\A(O)\Omega$. Consider $F(x_1,\dots,x_n)=\inner\Psi{\phi(x_1)\cdots\phi(x_n)\Omega}$, where translations $x_j-x_{j-1}$ are taken in $O$. By the spectrum condition (positive energy), $F$ extends to an analytic function of the differences $x_j-x_{j-1}$ in the forward tube $\mathrm{Im}\,(x_j-x_{j-1})\in V_+$. It vanishes on an open set (where all points are in $O$), hence everywhere by analyticity (edge-of-the-wedge/unique continuation). So $\Psi$ is orthogonal to the dense set generated by the vacuum, so $\Psi=0$.
\end{proofsketch}

Consequences. (i) \emph{Local operators in an arbitrarily small region can create any state}, including one with a particle at the other side of the universe; no contradiction with causality because the operator is not a local observable of that distant region and it needs a non-local looking combination of unitaries. (ii) No non-zero element of $\A(O)$ annihilates $\Omega$; so $\omega_0$ is a faithful state on $\A(O)$. (iii) Hence we may apply \cref{ch:tomita} to $(\A(O),\Omega)$ for \emph{every} region with non-trivial complement. (iv) Combined with \cref{ch:classification}: $\A(O)$ can have no minimal projection, since that would give a pure normal state incompatible with the analyticity underlying the proof; this is made precise in \cref{ch:typeIII_local}.

\section{Bisognano--Wichmann}

Let $W=\{x:x^1>|x^0|\}$ be the right Rindler wedge, $\A(W)$ its algebra, $\A(W')=\A(W)'$ (Haag duality) that of the left wedge. Let $\Lambda_W(s)$ be the boost with rapidity $s$ that preserves $W$ (moving forward in time inside $W$ for $s>0$), $U(\Lambda_W(s))=\ee^{\ii sB}$ with $B$ the boost generator.

\begin{theorem}[Bisognano--Wichmann]\label{thm:BW}
For the vacuum and a wedge algebra in a Wightman theory,
\begin{equation}
\Delta^{\ii t}_{\Omega}=U\big(\Lambda_W(-2\pi t)\big),\qquad J_\Omega=\Theta\,U(R_W(\pi)),
\end{equation}
where $\Theta$ is the CPT operator and $R_W(\pi)$ the rotation by $\pi$ in the plane orthogonal to the $(x^0,x^1)$ plane (so $J$ maps $\A(W)$ to $\A(W')$). In particular the modular Hamiltonian is $K=-\log\Delta=2\pi B$.
\end{theorem}
(This is a theorem for theories obeying the Wightman axioms; the result was proven in the 1970s by Bisognano and Wichmann, and extended to nets by various authors, so that $\Delta^{it}$ acts geometrically on the whole net of the wedge and its subregions.)

\begin{proofsketch}
Formally: the vacuum correlation functions are analytic in the forward tube. Rotating the Euclidean angle by $\pi$ in the $(x^0,x^1)$-plane maps the right wedge to the left wedge: $\Theta U(R(\pi))$ is the antilinear map $a\Omega\mapsto a\ad\Omega$, i.e.\ $S=J\Delta^{1/2}$ is the analytic continuation of the boost to ``imaginary rapidity $\ii\pi$'', so $\Delta^{1/2}=\ee^{-\pi B}$, which is $\ee^{-K/2}$ with $K=2\pi B$. Rigorous proofs use the spectrum condition to continue $U(\Lambda(s))$ analytically into the strip $0<\mathrm{Im}\,s<\pi$ on vectors $a\Omega$.
\end{proofsketch}

\begin{keyidea}
The modular flow of the Rindler wedge in the vacuum is the boost, with $K=2\pi B$ and $J$ the CRT reflection. The abstract Tomita--Takesaki objects $(\Delta,J)$ are Lorentz geometry.
\end{keyidea}

\section{The Unruh effect as modular flow}

Combining \cref{thm:BW} with the KMS property of \cref{ch:tomita}: the vacuum restricted to $\A(W)$ is KMS for the boost group at inverse temperature $\beta=2\pi$ (\cref{ch:kms}). An observer with uniform acceleration $a$ has trajectory $x(\tau)=(a^{-1}\sinh a\tau,a^{-1}\cosh a\tau)$ in $W$, and proper time $\tau$ corresponds to rapidity $s=a\tau$. The proper-time evolution is then $\alpha_\tau=\Ad U(\Lambda_W(a\tau))$, which is KMS at $\beta_\tau=2\pi/a$. A detector coupled to the field along this worldline therefore sees the vacuum as a thermal state at
\begin{equation}
T_U=\frac{a}{2\pi}\quad(\hbar=c=k_B=1).
\end{equation}
Here the temperature is a property of the vacuum's modular structure and not of any particle content: the vacuum \emph{is} an equilibrium state of the boost dynamics, in a representation (the vacuum Fock space) in which $\A(W)$ has no density matrix. The same logic gives the Hawking temperature as the KMS temperature at a Killing horizon (\cref{ch:curved_hawking}) and the Gibbons--Hawking temperature in the de Sitter static patch (\cref{ch:static_patch}).

\begin{whatwrong}
``The accelerated observer sees particles'' is a statement about the choice of the \emph{state restricted to the wedge}; the observer's algebra is $\A(W)$ in both descriptions. The inequivalence of the thermal Rindler and the Minkowski particle notions is the Fock-space inequivalence of \cref{ch:weyl}: Rindler quanta are not a unitary rotation of Minkowski quanta. Equally, ``the Unruh effect needs a horizon'' is built in: $\A(W)$ is a type III algebra \emph{because} the wedge has an entangling surface.
\end{whatwrong}

\section*{Exercises}
\begin{exercise}
For the massless scalar in $1+1$ dimensions, the right wedge algebra is generated by right- and left-moving chiral fields on half-lines. Show that on the one-particle space the boost acts as a dilation $p\mapsto\ee^{s}p$ of the null momentum and that $\Delta^{1/2}=\ee^{-\pi B}$ maps $p>0$ momenta as an analytic continuation to $\ee^{\ii\pi}p$.
\end{exercise}
\begin{exercise}
Using $K=2\pi B$ show that the vacuum expectation $\langle\Omega|K|\Omega\rangle=0$ while $\langle\Psi|K|\Psi\rangle>0$ for a one-particle excitation of positive boost energy located in $W$.
\end{exercise}
\begin{exercise}
Derive $T_U=a/2\pi$ from $\beta=2\pi$ in rapidity and $s=a\tau$. Restore $\hbar,c,k_B$.
\end{exercise}
\begin{exercise}
Argue from Reeh--Schlieder that the vacuum state restricted to $\A(O)$ is faithful and conclude it is not a vector state of a pure normal state: if it were, $\Omega$ would be a minimal-projection vector.
\end{exercise}
